\documentclass[10pt]{article}
\usepackage[utf8]{inputenc}
\usepackage[vietnamese]{babel}
\usepackage{amsmath,amssymb}
\usepackage{geometry}
\usepackage{graphicx}
\usepackage{tikz}
\usepackage{enumitem}
\usepackage{tabularx}
\usepackage{tcolorbox}
\tcbuselibrary{skins}

\geometry{
  paperwidth=8.5in,
  paperheight=11in,
  left=0.65in,
  right=0.65in,
  top=0.6in,
  bottom=0.55in
}

\definecolor{projgreen}{RGB}{0, 128, 105}
\definecolor{projbg}{RGB}{237, 246, 240}
\definecolor{iconblue}{RGB}{30, 130, 195}

\newcommand{\graphicon}{%
  \begin{tikzpicture}[baseline=-2pt, scale=0.85]
    \draw[iconblue, fill=iconblue!10, line width=0.6pt, rounded corners=0.8pt] (0,0) rectangle (9pt, 9pt);
    \draw[iconblue, line width=0.6pt] (1.5pt, 2pt) -- (3.5pt, 7pt) -- (5.5pt, 3.5pt) -- (7.5pt, 7.5pt);
  \end{tikzpicture}%
}

\begin{document}
\pagestyle{empty}
\fontsize{9.5pt}{12pt}\selectfont

% --- Header ---
\noindent
{\bfseries 260}\qquad{\bfseries\sffamily CHƯƠNG 3}\quad{\sffamily Các quy tắc đạo hàm}
\vspace{12pt}

% --- Main 2-Column Section (Asymmetric Layout) ---
\noindent
\begin{minipage}[t]{0.33\textwidth}
  \textbf{48.} Khi máu chảy dọc theo một mạch máu, lưu lượng $F$ (thể tích máu trên một đơn vị thời gian chảy qua một điểm cho trước) tỉ lệ với lũy thừa bậc bốn của bán kính $R$ của mạch máu:
  \[
    F = kR^4
  \]
  (Điều này được gọi là \emph{Định luật Poiseuille}; chúng ta sẽ giải thích tại sao định luật này đúng trong Mục 8.4.) Một động mạch bị tắc nghẽn một phần có thể được mở rộng bằng một thủ thuật gọi là nong mạch (\emph{angioplasty}), trong đó một ống thông có bóng ở đầu được bơm căng bên trong động mạch để làm rộng động mạch và phục hồi lưu lượng máu bình thường.
  \par\vspace{3pt}
  Hãy chỉ ra rằng biến thiên tương đối của $F$ xấp xỉ bằng bốn lần biến thiên tương đối của $R$. Việc tăng $5\%$ bán kính sẽ ảnh hưởng như thế nào đến lưu lượng máu?

  \vspace{10pt}
  \textbf{49.} Thiết lập các quy tắc sau khi làm việc với vi phân (trong đó $c$ ký hiệu một hằng số, $u$ và $v$ là các hàm số theo $x$).\vspace{3pt}

  \noindent
  \begin{tabularx}{\linewidth}{@{}X X@{}}
    (a) $dc = 0$ & (b) $d(cu) = c\,du$ \\[4pt]
    (c) $d(u + v) = du + dv$ & (d) $d(uv) = u\,dv + v\,du$ \\[4pt]
    (e) $d\left(\dfrac{u}{v}\right) = \dfrac{v\,du - u\,dv}{v^2}$ & (f) $d(x^n) = n x^{n-1}\,dx$
  \end{tabularx}

  \vspace{10pt}
  \textbf{50.} Trong các sách giáo trình vật lý, chu kỳ $T$ của một con lắc có chiều dài $L$ thường được cho bởi công thức $T \approx 2\pi\sqrt{L/g}$, với điều kiện con lắc dao động qua một cung tương đối nhỏ. Trong quá trình thiết lập công thức này, phương trình $a_T = -g\sin\theta$ cho gia tốc tiếp tuyến của quả nặng con lắc được rút ra, và sau đó $\sin\theta$ được thay thế bởi $\theta$ cùng với nhận xét rằng đối với các góc nhỏ, $\theta$ (tính bằng radian) rất gần với $\sin\theta$.
\end{minipage}%
\hfill
\begin{minipage}[t]{0.63\textwidth}
  \begin{enumerate}[label=\textbf{(\alph*)}, leftmargin=18pt, itemsep=4pt, topsep=0pt]
    \item Xác minh xấp xỉ tuyến tính tại 0 đối với hàm số sin:
    \[
      \sin\theta \approx \theta
    \]
    \item Nếu $\theta = \pi/18$ (tương đương với $10^\circ$) và chúng ta xấp xỉ $\sin\theta$ bởi $\theta$, thì sai số phần trăm là bao nhiêu?
    \item \graphicon\; Sử dụng đồ thị để xác định các giá trị của $\theta$ sao cho $\sin\theta$ và $\theta$ sai khác nhau ít hơn $2\%$. Các giá trị này tính theo độ là bao nhiêu?
  \end{enumerate}

  \vspace{8pt}
  \textbf{51.} Giả sử thông tin duy nhất mà chúng ta có về một hàm số $f$ là $f(1) = 5$ và đồ thị đạo hàm của nó như hình bên.
  \begin{enumerate}[label=\textbf{(\alph*)}, leftmargin=18pt, itemsep=2pt, topsep=2pt]
    \item Sử dụng xấp xỉ tuyến tính để ước lượng $f(0{,}9)$ và $f(1{,}1)$.
    \item Các ước lượng của bạn trong phần (a) là quá lớn hay quá nhỏ? Hãy giải thích.
  \end{enumerate}
  \vspace{2pt}
  \begin{center}
    \includegraphics[width=0.36\linewidth]{images/p0295_fig1.png}
  \end{center}

  \vspace{6pt}
  \textbf{52.} Giả sử chúng ta không có công thức của $g(x)$ nhưng biết rằng $g(2) = -4$ và $g'(x) = \sqrt{x^2 + 5}$ với mọi $x$.
  \begin{enumerate}[label=\textbf{(\alph*)}, leftmargin=18pt, itemsep=2pt, topsep=2pt]
    \item Sử dụng xấp xỉ tuyến tính để ước lượng $g(1{,}95)$ và $g(2{,}05)$.
    \item Các ước lượng của bạn trong phần (a) là quá lớn hay quá nhỏ? Hãy giải thích.
  \end{enumerate}
\end{minipage}

\vspace{16pt}

% --- DISCOVERY PROJECT SECTION ---
\noindent
\begin{tikzpicture}
  \node[anchor=west, inner sep=0pt] (t1) at (95pt, 0) {\color{projgreen}\sffamily\bfseries DỰ ÁN KHÁM PHÁ};
  \node[anchor=west, inner sep=0pt] (ico) at ([xshift=8pt]t1.east) {%
    \begin{tikzpicture}[baseline=-2pt]
      \draw[projgreen, line width=0.7pt, fill=white, rounded corners=0.8pt] (0,-1pt) rectangle (10pt, 9pt);
      \draw[projgreen, line width=0.5pt] (2pt, 1.5pt) -- (4.5pt, 6.5pt) -- (6.5pt, 3.5pt) -- (8.5pt, 7.5pt);
    \end{tikzpicture}%
  };
  \node[anchor=west, inner sep=0pt] (t2) at ([xshift=8pt]ico.east) {\color{projgreen}\sffamily\bfseries XẤP XỈ ĐA THỨC};
  \foreach \y in {-4.5, -1.5, 1.5, 4.5} {
    \draw[projgreen, line width=0.6pt] (0, \y pt) -- ([xshift=-10pt]t1.west |- 0, \y pt);
    \draw[projgreen, line width=0.6pt] ([xshift=10pt]t2.east |- 0, \y pt) -- (\linewidth, \y pt);
  }
\end{tikzpicture}

\vspace{-2pt}
\noindent
\begin{tcolorbox}[
  enhanced,
  colback=projbg,
  colframe=projbg,
  boxrule=0pt,
  arc=0pt,
  left=0pt,
  right=12pt,
  top=10pt,
  bottom=10pt,
  width=\linewidth
]
  \hspace*{0.33\linewidth}%
  \begin{minipage}{0.64\linewidth}
    \fontsize{9.2pt}{11.8pt}\selectfont
    Xấp xỉ bằng đường tiếp tuyến $L(x)$ là xấp xỉ bậc nhất (tuyến tính) tốt nhất cho $f(x)$ gần $x = a$ vì $f(x)$ và $L(x)$ có cùng tốc độ thay đổi (đạo hàm) tại $a$. Để có một xấp xỉ tốt hơn một đường thẳng tuyến tính, hãy thử một xấp xỉ bậc hai $P(x)$. Nói cách khác, chúng ta xấp xỉ một đường cong bằng một parabol thay vì một đường thẳng. Để đảm bảo rằng xấp xỉ là tốt, chúng ta đặt ra các điều kiện sau:
    
    \vspace{4pt}
    \begin{itemize}[leftmargin=*, labelsep=5pt, itemsep=2pt, topsep=2pt]
      \item[(i)] $P(a) = f(a)$ \hfill \makebox[220pt][l]{($P$ và $f$ phải có cùng giá trị tại $a$.)}
      \item[(ii)] $P'(a) = f'(a)$ \hfill \makebox[220pt][l]{($P$ và $f$ phải có cùng tốc độ thay đổi tại $a$.)}
      \item[(iii)] $P''(a) = f''(a)$ \hfill \makebox[220pt][l]{(Độ dốc của $P$ và $f$ phải thay đổi với cùng tốc độ tại $a$.)}
    \end{itemize}

    \vspace{4pt}
    \begin{enumerate}[label=\textbf{\arabic*.}, leftmargin=*, itemsep=5pt, topsep=4pt]
      \item Tìm xấp xỉ bậc hai $P(x) = A + Bx + Cx^2$ cho hàm số $f(x) = \cos x$ thỏa mãn các điều kiện (i), (ii), và (iii) với $a = 0$. Vẽ đồ thị của $P$, $f$, và xấp xỉ tuyến tính $L(x) = 1$ trên cùng một màn hình. Nhận xét về mức độ xấp xỉ $f$ của các hàm $P$ và $L$.
      \item Xác định các giá trị của $x$ sao cho xấp xỉ bậc hai $f(x) \approx P(x)$ trong Bài tập 1 chính xác trong phạm vi 0{,}1. [\textit{Gợi ý:} Vẽ đồ thị $y = P(x)$, $y = \cos x - 0{,}1$, và $y = \cos x + 0{,}1$ trên cùng một màn hình.]
      \item Để xấp xỉ một hàm số $f$ bằng một hàm bậc hai $P$ gần một số $a$, tốt nhất nên viết $P(x)$ dưới dạng
      \[
        P(x) = A + B(x - a) + C(x - a)^2
      \]
    \end{enumerate}
  \end{minipage}
\end{tcolorbox}

\vfill
\begin{center}
  \fontsize{7pt}{8pt}\selectfont\color{black!70}
  Copyright 2021 Cengage Learning. All Rights Reserved. May not be copied, scanned, or duplicated, in whole or in part. WCN 02-200-203
\end{center}

\end{document}